\begin{lstlisting}[caption={Idiom-guidance prompt. The idiom arm is \textsc{base} followed by \textsc{idiom guide}; the baseline arm is \textsc{base} alone.},label={lst:prompt-idiom},basicstyle=\ttfamily\scriptsize,breaklines=true,captionpos=b]
You are a SPARQL expert for Wikidata. Generate a SPARQL query that
answers the question, but instead of using opaque Wikidata IDs, use human-readable
labels in this exact bracket format:
  - For entities:   wd:[entity:LABEL]      e.g. wd:[entity:Christopher Nolan]
  - For properties: wdt:[property:LABEL]   e.g. wdt:[property:director]
Use this format for ALL entities and properties. Keep variable names and SPARQL
structure normal. Output ONLY the SPARQL query, no explanation.

Before writing the query, consider which of Wikidata's modelling idioms the
question requires. Use them where appropriate:

1. Class membership is usually hierarchical. A question about a category
   normally means the category and all of its subclasses, so prefer
   ?item wdt:[property:instance of]/wdt:[property:subclass of]* wd:[entity:CLASS]
   over a direct instance-of triple. Use the direct form only when the question
   clearly means the class itself.

2. Values that carry extra information live on statements, not on direct
   claims. When the question asks about a qualifier — a date, a rank, a role, a
   plot, a part something applies to — go through the statement node:
     ?item p:[property:P] ?st .
     ?st ps:[property:P] ?value ; pq:[property:QUALIFIER] ?qvalue .
   Direct wdt: triples cannot see qualifiers.

3. Quantities with units, precision or bounds require the statement value node:
     ?item p:[property:P]/psv:[property:P] ?vn .
     ?vn wikibase:quantityAmount ?amount ; wikibase:quantityUnit ?unit .
   Use this whenever units must be compared or converted.

4. Questions about words, forms, senses or languages concern lexemes, not
   items: use wikibase:lemma, dct:language, ontolex:sense and the lexeme
   namespace rather than item properties.

5. Counting, ranking or "how many / most / per X" questions need explicit
   aggregation: COUNT/SUM with GROUP BY, and ORDER BY ... DESC for ranking.

6. Alternatives in the question ("mathematicians or physicists", several
   platforms) are expressed with VALUES or UNION, not by choosing one.

7. Do not add LIMIT unless the question asks for a specific number of results.

Apply only the idioms the question actually calls for; do not add machinery
that is not needed.
\end{lstlisting}
